\section{Threats to Validity}
\label{sec:threats}
\textbf{Threats to construct validity} relate to the selection of jailbreak attacks and OOD thresholds.
We selected 10 different jailbreak attacks guided by their prevalence to guardrails and ability to represent a wide spectrum of jailbreak techniques \cite{dong2024safeguarding, yi2024jailbreak, wei2024jailbroken}. 
While additional attack types could be included in future evaluations, this would not fundamentally alter the key conclusions presented in our research questions: that complex jailbreak attacks can be effectively detected by our proposed OOD detection measures (RQ1), and that detected attacks can be efficiently learned by \ourapp through continual learning with minimal catastrophic forgetting (RQ2 and RQ3).
%
For the OOD threshold, we adopted commonly used quantile-based thresholds (90\%, 95\%, and 99\%) based on the distribution of OOD scores on the training set~\cite{liu2020energy,lee2018simple}. This approach allows us to select the most appropriate threshold values based on our experiments, ensuring that only the most anomalous prompts are flagged for guardrail adaptation. However, in practice, threshold selection is based on various factors depending on the application's tolerance for risk and the distributional properties of the data. 

\textbf{Threats to internal validity} relate to the potential influence of hyperparameter settings during the fine-tuning of our \ourapp~and the selection of OOD detection thresholds. Variations in energy threshold values, LoRA configurations, or learning rates, compared to those specified in Section \ref{sec:exp_design}, could impact experiment's outcomes. 
To address these threats, we open-source our replication package and provide detailed documentation of all hyperparameter settings to ensure the experiment is reproducible by future researchers. 

\textbf{Threats to external validity} concern the generalizability of our results. Our experiment findings are supported by the dataset, jailbreak methods, and model architecture employed during the study. Our training dataset contains 12,233 unsafe prompts and 6,905 safe prompts. Our  validation dataset contains 314 safe and 640 unsafe natural language prompts to assess false positive rates and catastrophic forgetting of~\ourapp. We also applied the 10 jailbreak attacks to 105 separate unsafe prompts, resulting in 1,050 jailbreak prompts across 10 categories for RQ1. Of the 1,050 jailbreak prompts, 992 are identified as OOD and used to study the continual learning adaptation in RQ2 and RQ3.
Our evaluation focuses on GPT-2 as the base model, which represents a widely-studied language model in AI safety research, but may not capture the behaviors of larger or more recent language models.
While \ourapp~is fine-tuned specifically to address the jailbreak attacks and model configurations discussed in this paper, other prompt datasets, jailbreak methods, and language model architectures can be explored in future work.


% \subsection{OOD Threshold Selection}
%  In our experiments, we adopted commonly used quantile-based thresholds (90\%, 95\%, and 99\%) based on the distribution of OOD scores on the training set~\cite{hendrycks2017baseline,liu2020energy,lee2018simple}. This approach allows us to select the most appropriate threshold values based on our experiments, ensuring that only the most anomalous prompts are flagged for guardrail adaptation. However, in practice, threshold selection are based on various factors depending on the application's tolerance for risk and the distributional properties of the data. Future works should focus on the different threshold selection methods to achieve the optimal performance on different metrics \cite{lipton2014optimal}. 